\begin{frame}{모델 기반이라는 전제: $P$ 를 ``정확히'' 안다}
  이번 장의 \textbf{모든 알고리즘}은 전이확률 $P(s'|s,a)$ 를
  \textbf{정확히 알고 있다는} 전제 위에 서 있다.
  \begin{block}{이 전제가 뜻하는 것}
    체스라면 \textbf{``이 수를 두면 상대는 어떻게 반응할
    확률이 얼마인가''를 미리 다 알아야} 한다는 뜻.
    \textbf{현실에서는 대부분 이걸 알 수 없다.}
  \end{block}
  \vskip0.5em
  \begin{itemize}
    \item \kb{Week 5$\cdot$6}은 이 \textbf{모델을 몰라도
      경험만으로} 학습하는 방법(몬테카를로, 시간차 학습).
    \item \kb{Week 9}는 ``상태가 너무 커서 표를 못 쓰는''
      문제를 신경망 근사로 다룬다.
  \end{itemize}
  \vskip0.5em
  \begin{block}{이 장을 받치는 뼈대}
    \textbf{최적해 = 벨만방정식의 고정점}. 모델 기반이라는
    \textbf{제약이 있지만}, 이후 모든 장의 \textbf{기준점}:
    Week 6 TD = 이 반복을 ``모델 대신 \textbf{경험으로}''
    근사한 것, Week 9 DQN = 같은 벨만 최적방정식의 고정점을
    \textbf{신경망}으로 찾아가는 것.
  \end{block}
\end{frame}
